\documentclass[reqno]{amsart}

\usepackage[T1]{fontenc}
\usepackage[utf8]{inputenc}
\usepackage{newtxtext,newtxmath}
\usepackage{microtype}
\usepackage{mathtools}
\usepackage{aliascnt}
\usepackage{enumitem}
\usepackage{booktabs,longtable,array}
\usepackage[numbers,sort&compress]{natbib}
\usepackage[hidelinks]{hyperref}
\usepackage[nameinlink,capitalise,noabbrev]{cleveref}
\makeatletter

% Journal-facing page style: no running heads, centered page numbers.
\pagestyle{plain}

% Preserve headline capitalization on the title page instead of AMS all-caps.
\def\@settitle{%
  \begin{center}%
    \baselineskip14\p@\relax
    \normalfont\Large\bfseries
    \@title\par
  \end{center}%
}

% Keep author names in normal case on the title page.
\def\@setauthors{%
  \begingroup
  \def\thanks{\protect\thanks@warning}%
  \trivlist
  \centering\footnotesize \@topsep26\p@\relax
  \advance\@topsep by -\baselineskip
  \item\relax
  \author@andify\authors
  \def\\{\protect\linebreak}%
  \authors
  \ifx\@empty\contribs
  \else
    ,\penalty-3 \space \@setcontribs
    \@closetoccontribs
  \fi
  \endtrivlist
  \endgroup
}

% Slightly calmer section spacing than the AMS default.
\def\section{\@startsection{section}{1}%
  \z@{2.2ex plus 0.8ex minus 0.2ex}{1.1ex}%
  {\normalfont\centering\scshape}}
\def\subsection{\@startsection{subsection}{2}%
  \z@{1.8ex plus 0.6ex minus 0.2ex}{0.8ex}%
  {\normalfont\bfseries}}
\def\subsubsection{\@startsection{subsubsection}{3}%
  \z@{1.5ex plus 0.4ex minus 0.2ex}{0.6ex}%
  {\normalfont\itshape}}

\makeatother

% Theorem styles with slightly more breathing room.
\newtheoremstyle{paperplain}%
  {7pt plus 2pt minus 1pt}% Space above
  {7pt plus 2pt minus 1pt}% Space below
  {\itshape}% Body font
  {}% Indent
  {\bfseries}% Head font
  {.}% Punctuation after head
  {0.6em}% Space after head
  {}% Head spec

\newtheoremstyle{paperdef}%
  {7pt plus 2pt minus 1pt}% Space above
  {7pt plus 2pt minus 1pt}% Space below
  {\normalfont}% Body font
  {}% Indent
  {\bfseries}% Head font
  {.}% Punctuation after head
  {0.6em}% Space after head
  {}% Head spec

\newtheoremstyle{paperremark}%
  {7pt plus 2pt minus 1pt}% Space above
  {7pt plus 2pt minus 1pt}% Space below
  {\normalfont}% Body font
  {}% Indent
  {\itshape}% Head font
  {.}% Punctuation after head
  {0.6em}% Space after head
  {}% Head spec

\theoremstyle{paperplain}
\newtheorem{theorem}{Theorem}[section]

\newaliascnt{proposition}{theorem}
\newtheorem{proposition}[proposition]{Proposition}
\aliascntresetthe{proposition}

\newaliascnt{lemma}{theorem}
\newtheorem{lemma}[lemma]{Lemma}
\aliascntresetthe{lemma}

\newaliascnt{corollary}{theorem}
\newtheorem{corollary}[corollary]{Corollary}
\aliascntresetthe{corollary}

\theoremstyle{paperdef}
\newaliascnt{definition}{theorem}
\newtheorem{definition}[definition]{Definition}
\aliascntresetthe{definition}

\newaliascnt{example}{theorem}
\newtheorem{example}[example]{Example}
\aliascntresetthe{example}

\theoremstyle{paperremark}
\newaliascnt{remark}{theorem}
\newtheorem{remark}[remark]{Remark}
\aliascntresetthe{remark}

\setlist[enumerate]{itemsep=2pt,topsep=4pt,parsep=0pt,partopsep=0pt}
\setlist[itemize]{itemsep=2pt,topsep=4pt,parsep=0pt,partopsep=0pt}

\setlength{\emergencystretch}{1.5em}
\allowdisplaybreaks[1]

\newcommand{\A}{\mathcal{A}}
\newcommand{\B}{\mathcal{B}}
\newcommand{\E}{\mathcal{E}}
\newcommand{\Reg}{\mathfrak{R}}
\newcommand{\dom}{\operatorname{dom}}
\newcommand{\Pow}{\mathcal{P}}

\crefname{theorem}{Theorem}{Theorems}
\Crefname{theorem}{Theorem}{Theorems}
\crefname{proposition}{Proposition}{Propositions}
\Crefname{proposition}{Proposition}{Propositions}
\crefname{lemma}{Lemma}{Lemmas}
\Crefname{lemma}{Lemma}{Lemmas}
\crefname{corollary}{Corollary}{Corollaries}
\Crefname{corollary}{Corollary}{Corollaries}
\crefname{definition}{Definition}{Definitions}
\Crefname{definition}{Definition}{Definitions}
\crefname{remark}{Remark}{Remarks}
\Crefname{remark}{Remark}{Remarks}
\crefname{example}{Example}{Examples}
\Crefname{example}{Example}{Examples}


\hypersetup{
  pdftitle={Fixed-Base Exhaustion for Internal Operators},
  pdfauthor={Anonymous},
  pdfsubject={Foundational logic manuscript},
  pdfkeywords={foundations of mathematics, conservative extension, external data, forcing, model theory, universal constructions}
}

\title[Fixed-Base Exhaustion]{Fixed-Base Exhaustion for Internal Operators}
\author[Anonymous]{Anonymous}
\date{}
\keywords{foundations of mathematics, conservative extension, external data, forcing, model theory, universal constructions}
\subjclass[2020]{Primary 03A05; Secondary 03C52, 03E40, 18A30}

\begin{document}

\begin{abstract}
We study fixed base structures equipped with an admissibility regime for partial operations and ask whether a carrier-preserving operator can yield genuinely new internal power on that same base. The paper proves a fixed-base exhaustion theorem: no carrier-preserving extension makes an old genuinely partial operation newly defined on old tuples, and once the effect of a purported operator is mediated by selectors, witnesses, evaluators, generic objects, or comparison data, the operator is conservative when those data are already fixed by the base and external otherwise. Applications cover canonical selection, completion, saturation and model-theoretic existence, forcing and absoluteness, and universal constructions.
\end{abstract}

\maketitle

\section{Introduction}\label{sec:introduction}

This paper asks a bounded foundational question: once a base structure is fixed, can a genuinely \emph{internal} operator create new admissible operator power on that same base? Here ``internal'' means carrier-preserving and not essentially dependent on witnesses, selectors, evaluators, or ambient objects imported from outside the original regime. The question arises because many standard methods---completion, saturation, forcing, universal constructions, and canonicalization---produce conclusions that are later read back to an original subject matter. The issue is not whether those methods are legitimate or fruitful. It is whether, relative to a fixed base, they amount to new internal operator power.

The central claim is negative and deliberately narrow. It is best read as a fixed-base exhaustion result rather than as a global anti-generative thesis. On a fixed base structure equipped with an admissibility regime, every purported carrier-preserving strengthening is either conservative on old-base-relevant content or else depends essentially on external data and so is not internal to that base. In particular, a carrier-preserving extension cannot totalize an operation that was genuinely partial on old tuples. The result is therefore a boundary theorem about internality, not a global thesis that external mathematics is dispensable.

The argument has three parts. First, in a partial-structure setting the paper proves a rigidity theorem: if an old primitive becomes newly defined on an old tuple, then the induced old structure has changed, so the move was not carrier-preserving after all. Second, it proves an auxiliary-data dichotomy: once the effect of a purported operator is mediated by a selector, witness, evaluator, completion boundary, or comparison map, one asks whether that datum is already fixed by the base. If it is, the operator is conservative; if not, the operator is external. Third, it proves an exhaustion theorem for the main kinds of apparent internal strengthening and then applies that scheme to several standard families.

The intended contrast is thus between internal generation and external method. Forcing works by adjoining generic objects and passing to a richer ambient universe; metric completion passes to a completed object; saturation and compactness work via larger models or witness structures; universal constructions are posed relative to ambient comparison data \citep{Hodges1993,Marker2002,Jech2003,Kanamori2008,MacLane1998}. None of this is challenged here. The claim is only that, measured against a fixed original base, such methods are external or scope-extending unless the final readback is conservative on the old base.

A terminological caution is useful. ``Admissibility'' is used here in a local, regime-relative sense, not as a synonym for admissibility in the sense of admissible set theory \citep{Barwise1975}. The comparison notion used throughout is simply conservativeness on old-base-relevant content: an operator or redescription adds no new values for old operations on old tuples and no new consequences about the old base beyond those already fixed by the original regime. In particular cases that notion may be sharpened to a more specific standard relation, such as definitional extension, conservative extension over the old signature, or interdefinability.

\paragraph*{Informal theorem.}
If a purported operator on a fixed base preserves the carrier and depends only on data already fixed by that base, then it is conservative on old-base-relevant content. If it requires genuinely external data, it is not internal to the base. In particular, no carrier-preserving extension of a nondegenerate partial base totalizes a genuinely partial old primitive on old tuples.

\paragraph*{Scope.}
The main text develops five flagship families: selection and canonical representatives, completion, model-theoretic saturation and existence, forcing and absoluteness, and existence by universal property. In each case the task is the same: isolate the load-bearing datum or ambient move and classify it as internally fixed, conservative readback, or explicit scope change. The paper does not claim that these families are mathematically shallow. It claims that, relative to a fixed base, their generative status is exhausted by that classification.

\paragraph*{Roadmap.}
\Cref{sec:prior-work} positions the paper relative to conservative extension, theory comparison, partial structures, model theory, forcing, universal constructions, and foundational methodology. \Cref{sec:formal-setting} gives the formal framework. \Cref{sec:rigidity} proves the no-internal-totalization theorem, and \Cref{sec:auxiliary-data} proves the auxiliary-data dichotomy. \Cref{sec:exhaustion} gives the exhaustion theorem for apparent strengthening mechanisms. \Cref{sec:case-studies} develops the flagship cases, \Cref{sec:no-emergence} shows that recombination does not generate a new family, and \Cref{sec:main-theorem} states the final synthesis. Appendix~\ref{app:extended-cases} records longer case analyses and mixed-family examples.

\section{Relation to Prior Work and Positioning}\label{sec:prior-work}

This paper sits at the junction of several established literatures. Its main theorem is new in its exact form, but its proof language and case studies draw on standard work in logic and foundations. The point of this section is therefore not to claim a single narrow ancestry. It is to make clear which literatures the manuscript enters and what it takes from each of them.

\paragraph*{Conservative extension and theory comparison.}
One strand of the paper concerns when an apparent strengthening says nothing genuinely new about the original base. That issue belongs to the familiar study of definitions, expansions, interpretations, and conservative extensions in model theory and adjacent work on theory comparison \citep{ChangKeisler1990,Hodges1993}. More recent discussions of definitional, categorical, and Morita equivalence sharpen the idea that two presentations may differ in signature while agreeing on what is added about their original subject matter \citep{BarrettHalvorson2016}. The present manuscript does not collapse those notions into one equivalence relation. It uses them as anchors for the narrower claim that internally fixed auxiliary data yield only conservative change.

\paragraph*{Partial structures and genuine undefinedness.}
The no-totalization result belongs formally to the setting of partial structures and partial algebras, where undefinedness is part of the structure rather than a temporary defect. Burmeister provides the relevant background vocabulary \citep{Burmeister1986}. The present use of partiality is sharply focused: it makes precise the claim that if an old primitive acquires a new value on an old tuple, then the old base itself has changed.

\paragraph*{External methods in model theory and set theory.}
Several of the paper's main examples are standard external methods of logic. Model theory uses compactness, saturation, monster models, and witness structures to move to richer environments \citep{Hodges1993,Marker2002}. Set theory uses forcing, generic extensions, and absoluteness to compare what survives passage to new universes \citep{Jech2003,Kanamori2008}. The paper does not dispute those practices. It reclassifies their fixed-base status: relative to a fixed base, they are ambient enlargements, witness imports, or conservative readbacks, not autonomous sources of new internal operator power.

\paragraph*{Universal constructions and ambient comparison.}
Category-theoretic existence by universal property is another natural neighboring literature. Universal arrows, adjunctions, limits, colimits, and free constructions are posed relative to an ambient category and comparison maps \citep{MacLane1998}. For the present paper, that ambient dependence is not an objection; it is the feature that marks the boundary between internal generation on an object alone and existence mediated by a wider framework.

\paragraph*{Foundational methodology.}
Methodologically, the closest analogue is fixed-base analysis in foundations. Reverse mathematics asks which set-existence principles are needed over a chosen base theory to prove a theorem \citep{Simpson2009}. The present paper asks a different but structurally similar question: once a base structure is fixed, which operator moves are already internal to it, and which rely on additional witnesses or scope extension? In that sense the manuscript is best read as a boundary theorem about internality rather than as an elimination program.

\paragraph*{Terminological caution and novelty.}
The manuscript's use of ``admissibility'' is local and regime-relative; it is not meant to place the paper inside admissible set theory unless that connection is made explicitly for some later application \citep{Barwise1975}. The novelty claim is likewise specific. The paper is not merely putting conservative extension, forcing, saturation, and universal constructions into one list. It argues that these apparently different families are governed by a common dichotomy: once the load-bearing datum is isolated, the move is either conservative because that datum was already fixed by the base, or external because it was not. Coupled with the rigidity theorem for partial structures and the exhaustion theorem for apparent strengthening mechanisms, that is the unified contribution of the manuscript.

\section{Formal Setting: Base Structures, Admissibility, and Internal Operators}\label{sec:formal-setting}

This section fixes the formal framework used throughout. The definitions are phrased in standard mathematical language while preserving the distinctions that drive the later proofs.

\begin{definition}[Base structure]
A \emph{base structure} is a triple
\[
\B=(X,\Omega,\A),
\]
where $X$ is a carrier, $\Omega$ is a family of partial operations
\[
\omega:X^{n_\omega}\rightharpoonup X
\]
of finite arities, and $\A$ is an admissibility regime specifying which primitive applications and compositions count as licensed on the carrier under discussion.
\end{definition}

\begin{remark}
The admissibility regime may be presented in more than one equivalent way: by primitive domains of definition, by closure constraints on compositions, or by an explicit class of licensed constructions. The essential point is that admissibility is part of the structure under analysis, not a temporary gap to be repaired later.
\end{remark}

\begin{definition}[Carrier-preserving operator]
Let $\B=(X,\Omega,\A)$ be a base structure. A purported operator $F$ on admissible inputs from $X$ is \emph{carrier-preserving} if it is presented as acting on the same base carrier $X$ rather than by openly replacing $X$ with a quotient, completion, extension, or new ambient structure.
\end{definition}

\begin{definition}[Conservative Action on Old-Base-Relevant Content]
A purported operator or redescription on a base structure $\B=(X,\Omega,\A)$ \emph{acts conservatively on old-base-relevant content} if it adds no new values for old operations on old tuples and no new consequences about the old base beyond those already fixed by the original regime.
\end{definition}

\begin{remark}
Depending on the application, conservative action may be sharpened into a more specific standard notion, such as definitional extension, interdefinability, or conservative extension over an old signature. The arguments below use only the common core of those notions: no new old-tuple values for old operations and no new old-base consequences.
\end{remark}

\begin{definition}[Internally fixed auxiliary datum]
Let $x\in X$ be an admissible input. An auxiliary datum $a(x)$ is \emph{internally fixed over $x$} if every admissibility-preserving redescription of the base structure that leaves the standing content of $x$ unchanged also leaves $a(x)$ unchanged up to the normalizations already licensed by the regime.
\end{definition}

\begin{definition}[External parameter or witness]
An auxiliary datum relevant to a purported operator on $x$ is \emph{external} if it is not internally fixed over $x$ and must therefore be supplied by an extra selector, witness, evaluator, generic object, comparison map, completion rule, or ambient construction not already determined by the original base structure.
\end{definition}

\begin{proposition}[Internality firewall]\label{prop:firewall}
Throughout this paper, the following principles are fixed.
\begin{enumerate}[label=(\roman*)]
    \item Admissibility is structural rather than merely provisional.
    \item A carrier-preserving redescription cannot create new values for old operations on old tuples.
    \item If a purported carrier-preserving operator depends essentially on external parameters or witnesses, then it is not internal to the original base structure.
\end{enumerate}
\end{proposition}

\begin{proof}
Clause (i) is built into the role assigned to the admissibility regime. Clause (ii) is exactly the content of conservative action on old-base-relevant content in this setting. Clause (iii) then follows immediately: if the effect of an operator depends on data not already fixed by the base, the operator's action is no longer determined by the base alone and so is not internal to it.
\end{proof}

\begin{remark}[Terminological boundary]
The manuscript's use of ``admissibility'' should not be read as a claim about admissible set theory in the sense of Barwise \citep{Barwise1975}. The intended notion here is regime-relative admissibility internal to a chosen base structure.
\end{remark}

\section{Rigidity of Carrier-Preserving Extension}\label{sec:rigidity}

We now formulate the first theorem engine in standard language. The underlying idea is simple: if an old partial operation becomes newly defined on an old tuple, then the old base structure has changed. The point of the section is to make that observation formal enough to support the later foundational claims.

\begin{definition}[Admissible extension]
Let $\B=(X,\Omega,\A)$ be a base structure. An \emph{admissible extension} of $\B$ is a structure
\[
\B^*=(X^*,\Omega^*,\A^*)
\]
such that:
\begin{enumerate}[label=(\roman*)]
    \item $X\subseteq X^*$;
    \item for each primitive $\omega\in\Omega$ there is a primitive $\omega^*\in\Omega^*$ of the same arity whose restriction to old defined tuples agrees with $\omega$;
    \item compositions of old primitives in the extension restrict to the old compositions on old admissible tuples; and
    \item the old admissibility profile is preserved as the induced profile on the old carrier unless the extension is explicitly announced as a change of carrier or regime.
\end{enumerate}
\end{definition}

\begin{definition}[Carrier-preserving extension]
An admissible extension $\B^*$ of $\B$ is \emph{carrier-preserving} if it is presented as leaving $X$ fixed as the base carrier under discussion rather than replacing $X$ by a quotient, completion, or new ambient domain.
\end{definition}

\begin{definition}[Nondegeneracy]
A base structure $\B=(X,\Omega,\A)$ is \emph{nondegenerate} if at least one primitive operation in $\Omega$ has a proper, non-total domain on $X$.
\end{definition}

\begin{lemma}[Old-graph rigidity]\label{lem:old-graph-rigidity}
Let $\B=(X,\Omega,\A)$ be a base structure and let $\B^*=(X^*,\Omega^*,\A^*)$ be an admissible extension. Suppose there exist a primitive $\omega\in\Omega$ of arity $n$ and a tuple $a\in X^n$ such that
\[
a\notin \dom(\omega)
\qquad\text{but}\qquad
 a\in \dom(\omega^*).
\]
Then the induced partial structure on the old carrier $X$ has changed. In particular, the move is not conservative on the old base.
\end{lemma}

\begin{proof}
The graph of an old primitive operation on old tuples is part of the old structure. If a tuple that was previously outside the domain of $\omega$ becomes newly defined under $\omega^*$, then the graph of that primitive has been enlarged on old data. That is not a mere change of notation or presentation. It is a change in the induced structure on $X$ itself.
\end{proof}

\begin{lemma}[Domain fixity on the old carrier]\label{lem:domain-fixity}
Let $\B=(X,\Omega,\A)$ be a nondegenerate base structure, and let $\B^*$ be a carrier-preserving extension of $\B$. Then for every primitive $\omega\in\Omega$ of arity $n$,
\[
\dom(\omega^*)\cap X^n = \dom(\omega).
\]
\end{lemma}

\begin{proof}
Because each old primitive is preserved by restriction, every old defined tuple remains defined, so
\[
\dom(\omega)\subseteq \dom(\omega^*)\cap X^n.
\]
Assume the inclusion is strict. Then some old tuple becomes newly defined in the extension. By \Cref{lem:old-graph-rigidity}, the induced partial structure on the old carrier has changed. But a carrier-preserving extension is, by definition, one that leaves the old carrier and its old induced structure in place rather than replacing it by a new regime. So strict inclusion is impossible.
\end{proof}

\begin{theorem}[No internal totalization]\label{thm:no-internal-totalization}
Let $\B=(X,\Omega,\A)$ be a nondegenerate base structure. No carrier-preserving extension of $\B$ totalizes a genuinely partial old operation on old tuples.
\end{theorem}

\begin{proof}
Suppose some old primitive $\omega\in\Omega$ is partial on $X$ but becomes total on old tuples in a carrier-preserving extension. Then
\[
\dom(\omega^*)\cap X^n = X^n
\qquad\text{while}\qquad
\dom(\omega)\subsetneq X^n,
\]
contradicting \Cref{lem:domain-fixity}.
\end{proof}

\begin{corollary}[No post hoc repair]\label{cor:no-posthoc}
An old undefined primitive application cannot be made admissible later by a carrier-preserving extension of the same base structure.
\end{corollary}

\begin{proof}
Any such repair would assign a new value to an old tuple for an old primitive operation, which is the graph change excluded by \Cref{lem:domain-fixity,thm:no-internal-totalization}.
\end{proof}

\begin{corollary}[No deferred admissibility]\label{cor:no-deferred}
If an old tuple is outside the domain of an old operation at the stage of formation, no later carrier-preserving extension makes that old application defined on the same old data.
\end{corollary}

\begin{proof}
Deferral changes only the timing of the same forbidden move. The old tuple would still have to become newly defined on the old carrier.
\end{proof}

\begin{remark}
This section provides the first formal reason the broader foundational thesis is plausible. Any later construction that appears to rescue an old undefined case must either fail, collapse to conservative action on old-base-relevant content, or else reveal itself as a move to a different carrier or a richer ambient setting.
\end{remark}

\section{The Internal/External Auxiliary-Data Theorem}\label{sec:auxiliary-data}

The rigidity theorem handles direct attempts to change the values of old operations on old tuples. A second and more conceptual issue remains. Many purported internal operators do not advertise themselves as changing an old graph. Instead they proceed by introducing an auxiliary selector, witness, evaluator, completion rule, comparison map, or generic object. The right question is then: was that auxiliary datum already fixed by the original base, or has it been supplied from outside?

\begin{theorem}[Auxiliary-data dichotomy]\label{thm:auxiliary-dichotomy}
Let $F$ be a carrier-preserving operator on admissible input $x\in X$ for a base structure $\B=(X,\Omega,\A)$. Suppose that the old-base-relevant content of $F(x)$ is determined, up to conservative action on old-base-relevant content, by $x$ together with an auxiliary datum $a(x)$. Then:
\begin{enumerate}[label=(\alph*)]
    \item if $a(x)$ is internally fixed over $x$, then $F$ acts conservatively on the old-base-relevant content of $x$;
    \item if $a(x)$ is external, then $F$ is not internal to the base structure.
\end{enumerate}
\end{theorem}

\begin{proof}
Assume first that $a(x)$ is internally fixed over $x$. Then every admissibility-preserving redescription of the base that leaves the relevant content of $x$ unchanged also leaves $a(x)$ unchanged up to licensed normalization. So $a(x)$ contributes no genuinely new distinction beyond what was already fixed by the base together with $x$. The effect of $F$ is therefore only to reorganize or read off structure that was already there, which is to say $F$ is conservative on the old base.

Assume instead that $a(x)$ is external. Then the effect of $F$ depends essentially on a datum not determined by the old base together with $x$. By Proposition~\ref{prop:firewall}, such dependence is incompatible with internality relative to the original base structure. Hence $F$ is not internal to that base.
\end{proof}

\begin{corollary}[No internalization by relabeling]\label{cor:no-internalization-by-relabeling}
If an external datum is redescribed as latent, semantic, structural, or canonical without becoming internally fixed over the original input, its status does not change. Any operator using it remains external to the base.
\end{corollary}

\begin{proof}
A change of description does not make an unfixed datum fixed. If the relevant effect still depends on data not determined by the original base, then the operator still depends on an external parameter or witness.
\end{proof}

\begin{remark}
This theorem is the conceptual center of the manuscript. Every case study in the later sections is meant to answer the same question in a different setting: which auxiliary datum does the purported construction need, and is that datum already fixed by the original base or imported from outside?
\end{remark}

\begin{example}[Canonicalization]
Suppose an operator assigns to each equivalence class a distinguished representative. If the representative rule is already fixed by the base structure, the operator is conservative. If the rule must be supplied from outside, then the operator is external by \Cref{thm:auxiliary-dichotomy}.
\end{example}

\begin{example}[Completion]
Suppose an operator assigns to an object its completed version together with a canonical embedding. If the completion is genuinely a new ambient object or depends on an external comparison rule, then the move is not internal to the original base. The later case-study section will make this more precise.
\end{example}

\section{Exhaustion of Apparent Internal Extension Mechanisms}\label{sec:exhaustion}

\Cref{thm:no-internal-totalization} blocks direct totalization of old partial operations, and \Cref{thm:auxiliary-dichotomy} classifies operators by the status of the data they require. The next task is to show that the main kinds of apparent internal strengthening one might try are exhausted by a short list of mechanisms.

\begin{definition}[Five extension mechanisms]\label{def:five-mechanisms}
A purported strategy for obtaining new internal operator power on a fixed base typically proceeds, if it succeeds at all, by one or more of the following moves.
\begin{enumerate}[label=\textup{(\Roman*)}, leftmargin=3em]
    \item \emph{Domain extension on old tuples}: an old primitive acquires new values on old inputs.
    \item \emph{Carrier change or quotient}: the effective object under discussion becomes a quotient, completion, extension, or otherwise different carrier.
    \item \emph{Evaluative closure or completion}: an added global evaluator, comparison rule, semantic predicate, or completion device decides cases that were previously not internally determined.
    \item \emph{Conservative definitional expansion}: the language changes, but the old base and its old consequences do not.
    \item \emph{New operator family or admissibility policy}: the regime itself is enlarged in a way not reducible to the first four cases.
\end{enumerate}
\end{definition}

\begin{theorem}[Exhaustion theorem]\label{thm:exhaustion-theorem}
Every apparent internal strengthening of a fixed base structure falls under at least one of the five mechanisms in \Cref{def:five-mechanisms}.
\end{theorem}

\begin{proof}
Let $\tau$ be a purported strengthening. If it makes an old primitive newly defined on an old tuple, it is of type (I). If the appearance of strength comes only because the object under discussion has been replaced by a quotient, completion, extension, or other new carrier, it is of type (II). If previously unresolved cases are decided by an added evaluator, global comparison, semantic oracle, or completion procedure, it is of type (III). If none of those changes occurs and only the descriptive apparatus changes, it is of type (IV). Any remaining case changes the operative family of constructions or the admissibility regime itself, which is type (V).
\end{proof}

\begin{remark}[Coverage rather than partition]
The five mechanisms are not claimed to be mutually exclusive. A single construction may instantiate more than one mechanism. Exhaustive coverage is enough for the argument.
\end{remark}

\begin{proposition}[Disposition of the five mechanisms]\label{prop:five-disposition}
Relative to the framework of \Cref{sec:formal-setting}, each mechanism in \Cref{def:five-mechanisms} is either blocked or else reveals that no genuinely internal strengthening has occurred.
\begin{enumerate}[label=\textup{(\Roman*)}, leftmargin=3em]
    \item Domain extension on old tuples is impossible by \Cref{thm:no-internal-totalization}.
    \item Carrier change or quotient is, by description, not carrier-preserving.
    \item Evaluative closure or completion depends on additional data or evaluators and is therefore external unless it collapses to conservative action on old-base-relevant content.
    \item Conservative definitional expansion is explicitly conservative and adds no new internal power.
    \item A new operator family or admissibility policy is an open change of regime rather than an internal strengthening of the original base.
\end{enumerate}
\end{proposition}

\begin{proof}
Clause (I) is exactly the rigidity theorem. Clause (II) restates the meaning of carrier preservation. Clause (III) is an application of the auxiliary-data dichotomy: if the evaluator or completion rule is already fixed by the base, the resulting move is conservative; if not, the move is external. Clause (IV) follows from the definition of conservative action on old-base-relevant content. Clause (V) is a change of the very regime relative to which internality was being measured.
\end{proof}

\begin{corollary}[Closure of internal totalization]\label{cor:closure-totalization}
No apparent internal strengthening survives as a genuine carrier-preserving totalization of old partial operations on old tuples.
\end{corollary}

\begin{proof}
Any such strengthening would fall under one of the five mechanisms by \Cref{thm:exhaustion-theorem}. By Proposition~\ref{prop:five-disposition}, each mechanism is either blocked or else does not yield genuine internal strengthening.
\end{proof}

\begin{remark}
At this point the manuscript has its basic formal spine. The remaining work is to instantiate that spine in the main case studies, then prove that combinations of the case-study mechanisms do not create a new internal family by interaction.
\end{remark}

\section{Flagship Case Studies}\label{sec:case-studies}

With the formal spine in place, we can now test the framework on five familiar families. The aim is not to rehearse their internal mathematics in full. It is to identify, for each family, the load-bearing datum or ambient move from the standpoint of a fixed base structure. Once that step is isolated, \Cref{thm:no-internal-totalization} and \Cref{thm:auxiliary-dichotomy}, together with Proposition~\ref{prop:five-disposition}, do the substantive work.

\begingroup
\small
\setlength{\LTleft}{0pt}
\setlength{\LTright}{0pt}
\renewcommand{\arraystretch}{1.12}
\begin{longtable}{>{\raggedright\arraybackslash}p{0.23\textwidth} >{\raggedright\arraybackslash}p{0.34\textwidth} >{\raggedright\arraybackslash}p{0.27\textwidth}}
\caption{Fixed-base classification of the flagship families.}\label{tab:flagship-families}\\
\toprule
Family & Load-bearing datum or move & Fixed-base classification \\
\midrule
\endfirsthead
\multicolumn{3}{@{}l}{\footnotesize\itshape Table \thetable{} continued.}\\
\toprule
Family & Load-bearing datum or move & Fixed-base classification \\
\midrule
\endhead
Selection, canonical representatives, quotienting & Representative rule, normalization criterion, or quotient passage & Conservative only if the rule is internally fixed; otherwise external. Quotient passage changes carrier. \\
Completion and evaluative closure & Completion object, embedding, boundary, or evaluator & Either external by carrier change, or conservative on old-base-relevant content if the evaluative rule is already fixed. \\
Compactness, saturation, model-theoretic existence & Larger model, realized type, witness structure, or existence theorem & Conservative only in the vacuous/definable case; otherwise external ambient enlargement or witness import. \\
Forcing, genericity, absoluteness & Generic object, forcing extension, or transfer criterion & Generic output is external; absoluteness/read-back is conservative on old-base-relevant content. \\
Universal constructions & Ambient category, comparison maps, universal witness object & Conservative only if the relevant witness is already fixed; otherwise ambiently external and typically carrier-changing. \\
\bottomrule
\end{longtable}
\endgroup

The subsections justify the last column of the table.

\subsection{Selection, canonical representatives, and quotienting}

Selection and canonicalization are the cleanest test cases for the auxiliary-data theorem. A base structure may determine an equivalence relation, a normalization target, or a family of interchangeable witnesses. The foundational question is whether it also determines a distinguished representative, or whether that representative is chosen by an additional ranking, coding, ordering, or normalization rule.

\begin{proposition}[Canonicalization, selection, and quotienting]\label{prop:case-canonicalization}
Let $D\subseteq X$ be an admissible domain for a base structure $\B=(X,\Omega,\A)$, and let $\sim$ be an equivalence relation on $D$ invariant under admissibility-preserving redescription. Suppose a purported carrier-preserving operator $N:D\to D$ satisfies
\[
N(x)\sim x \qquad\text{and}\qquad x\sim y \Rightarrow N(x)=N(y).
\]
Then the following trichotomy governs the case.
\begin{enumerate}[label=(\alph*)]
    \item If the representative rule determining $N(x)$ is internally fixed over $x$, then $N$ is conservative on the old-base-relevant content of $x$.
    \item If that rule is not internally fixed, then $N$ depends on an external selector, ranking, or normalization datum and is not internal to $\B$.
    \item If the maneuver is really passage from $D$ to the quotient $D/{\sim}$, then the carrier has changed and the move is not carrier-preserving in the first place.
\end{enumerate}
\end{proposition}

\begin{proof}
If $N$ is genuinely a map from $D$ to $D$, then its effect is mediated by the rule that chooses, for each equivalence class, which representative counts as canonical. That rule is exactly the auxiliary datum relevant to the construction. If the datum is internally fixed, then \Cref{thm:auxiliary-dichotomy} yields conservativity: the operator merely reports a representative already determined by the base. If the datum is not internally fixed, then the choice of representative depends on an external ranking, code, or normalization criterion, so \Cref{thm:auxiliary-dichotomy} makes the operator external to $\B$.

If, instead, the supposed canonicalization works by replacing each element with its equivalence class, then the effective object under discussion is no longer $D$ but the quotient $D/{\sim}$. That is a carrier change and so falls under mechanism (II) from \Cref{def:five-mechanisms}, not under internal strengthening of the original base.
\end{proof}

\begin{example}[Least representative rules]
Choosing ``the least code,'' ``the lexicographically first form,'' or ``the minimal-complexity representative'' is internal only relative to a base that already fixes the relevant order, coding, or complexity measure. Otherwise the apparent canonicity is carried by imported structure rather than by the original base itself.
\end{example}

\subsection{Completion and evaluative closure}

Completion constructions are especially instructive because they often feel canonical while still operating by passage to a richer ambient object. In metric, order-theoretic, algebraic, and semantic settings, completion typically works either by adjoining new points, classes, or ideal objects, or by adding a rule that evaluates old data against a richer completion boundary.

\begin{proposition}[Completion and evaluative closure]\label{prop:case-completion}
Let $C$ be a purported completion operator applied to an admissible object $A$ over a fixed base structure. Then exactly one of the following happens.
\begin{enumerate}[label=(\alph*)]
    \item $C(A)$ is genuinely a richer object---for example, a completion, closure object, or ideal enlargement with new points, classes, or totalizations---in which case the move is not carrier-preserving.
    \item The only old-base-relevant effect of the completion is mediated by an internally fixed comparison, embedding, or evaluative rule, in which case the move is conservative on the old base.
    \item The read-back from the completion depends on a comparison, embedding, boundary, or evaluator not internally fixed by the original base, in which case the move is external.
\end{enumerate}
\end{proposition}

\begin{proof}
If the completion produces a genuinely new object, then the carrier under discussion has changed. A metric completion, a Dedekind or Cauchy-style enlargement, a closure by ideal elements, or any analogous passage to a richer ambient object is therefore external from the standpoint of the original base. This is mechanism (II) from \Cref{def:five-mechanisms}.

Suppose instead that the completion is used only to obtain old-base-relevant consequences: for example, to justify a limit statement about old parameters, to extend a comparison already implicit in the base, or to read back a property of the original object. Then the effect is governed by an auxiliary datum consisting of the relevant embedding, evaluator, or completion boundary. If that datum is internally fixed, \Cref{thm:auxiliary-dichotomy} shows that the move is conservative on the old base. If it is not internally fixed, the completion works only by importing exactly the additional evaluative structure on which the result depends, and the move is external.
\end{proof}

\begin{example}[Metric completion]
The passage from $\mathbb{Q}$ to $\mathbb{R}$ is mathematically canonical, but it is still a passage to a richer object rather than an internal operator on the original carrier $\mathbb{Q}$. Statements about rationals that are later justified by working in $\mathbb{R}$ may be highly informative, but from the fixed-base standpoint they are either conservative read-backs or consequences of an external ambient enlargement.
\end{example}

\begin{remark}
Uniqueness up to isomorphism or isometry does not alter this diagnosis. It controls ambiguity in the richer ambient object; it does not turn that richer object into a same-carrier operator on the original base.
\end{remark}

\subsection{Compactness, saturation, and model-theoretic existence}

Model theory supplies many of the clearest examples of mathematically fruitful external method: compactness arguments, saturated and monster models, elementary extensions, and existence proofs by witness structures \citep{ChangKeisler1990,Hodges1993,Marker2002}. The present claim is not that those methods are dispensable. It is that, from the standpoint of a fixed base structure, they do not generate new internal operator power unless the relevant witness was already fixed.

\begin{proposition}[Compactness, saturation, and model-theoretic existence]\label{prop:case-model-theory}
Let $M$ be a structure, and let $F$ be a purported operator whose effect is justified by compactness, saturation, a monster-model argument, or a model-existence theorem. Then one of the following holds.
\begin{enumerate}[label=(\alph*)]
    \item The relevant witness, realization, or map was already definably fixed on the old base, in which case $F$ is conservative on the old-base-relevant content of $M$.
    \item The construction proceeds by passing to a larger model, adjoining new realizing elements, invoking a witness structure, or using an existence theorem whose witness is not part of $M$, in which case $F$ is external to the original base.
\end{enumerate}
In particular, none of these model-theoretic methods yields new internal operator power on a fixed old carrier except in the vacuous conservative case.
\end{proposition}

\begin{proof}
Consider first type realization and saturation. If the purported operator produces a realization already available in $M$, then nothing new has been generated; the effect is conservative. If the realization is not already present, then some new witness element or larger ambient model is required. That is a change of carrier or witness import, not an internal operator on the original base.

The same point applies to saturated envelopes, monster models, and compactness arguments. A monster model is explicitly a richer ambient structure. A compactness argument that establishes the existence of some model or element does so by producing or appealing to a witness structure not identical with the original base. If the only conclusion read back to $M$ is a property already fixed by the old base, then the use of the external witness is conservative on old-base-relevant content. Otherwise the load-bearing datum is external, and \Cref{thm:auxiliary-dichotomy} classifies the move accordingly.
\end{proof}

\begin{example}[Monster models]
Working inside a monster model can make independence, type comparison, and canonical parameter arguments much cleaner. But the monster model is an ambient enlargement. It is not an internal operator on the original small structure merely because later conclusions are stated back in the original language.
\end{example}

\subsection{Forcing, genericity, and absoluteness}

Forcing is the clearest case in which a foundational paper must be careful not to overclaim. The method is one of the great successes of modern logic, and much of its power comes precisely from controlled passage to a richer ambient universe \citep{Jech2003,Kanamori2008}. The fixed-base diagnosis is correspondingly sharp: a forcing extension is an external move, while read-back by absoluteness or invariance is conservative on the old base rather than internally generative.

\begin{proposition}[Forcing, genericity, and absoluteness]\label{prop:case-forcing}
Let $M$ be a ground model or other fixed set-theoretic base, and let $F$ be a purported forcing-based operator.
\begin{enumerate}[label=(\alph*)]
    \item If $F$ works by adjoining a generic filter, generic real, or forcing extension $M[G]$, then it is external to $M$ by imported witness or ambient enlargement.
    \item If $F$ uses forcing only to transfer a statement about old parameters by absoluteness, invariance, or ground-model forcing analysis, then its effect on old-base-relevant content is conservative rather than a new internal operator on $M$.
\end{enumerate}
Hence forcing and genericity do not provide new internal operator power on the original base model.
\end{proposition}

\begin{proof}
In the first case the point is immediate. A generic object $G$ is not part of the original base, and the passage from $M$ to $M[G]$ is an explicit move to a richer ambient universe. That is the paradigmatic instance of an external method.

In the second case the conclusion concerns only old parameters or old-base-relevant statements. Then forcing functions as a proof method or transfer method: one studies what is preserved, reflected, or already decided when the situation is viewed through forcing. The result is therefore not a new internal operator on $M$, but a conservative read-back about the old base. Boolean-valued, algebraic, or completion-style presentations of forcing do not change this point. They either redescribe already fixed information conservatively or shift the work onto additional evaluative or completion structure, which is external by Proposition~\ref{prop:five-disposition} and Proposition~\ref{prop:case-completion}.
\end{proof}

\begin{example}[Generic objects versus absolute statements]
A Cohen real added over a ground model is not an internal object of that ground model. By contrast, an absoluteness argument that transfers a statement about old natural numbers back to the ground model does not thereby create a new internal operator on the ground model; it certifies old-base-relevant content by an external route.
\end{example}

\subsection{Universal constructions and existence by universal property}

Universal constructions are often described as canonical, and they are canonical in an important categorical sense. But their canonicity is mediated by an ambient category, a class of comparison maps, and a universal condition imposed across that ambient framework \citep{MacLane1998}. The fixed-base question is therefore not whether the universal construction is mathematically legitimate, but whether the relevant universal witness is already fixed by the old base or supplied through ambient categorical data.

\begin{proposition}[Universal constructions and existence by universal property]\label{prop:case-universal}
Let $A$ be an object presented by a fixed base structure, and suppose a purported operator assigns to $A$ an object $U(A)$ by a universal property in an ambient category $\mathcal{C}$. Then the following distinction is exhaustive.
\begin{enumerate}[label=(\alph*)]
    \item If the universal witness and its comparison maps are already fixed, relative to the chosen ambient framework, strongly enough that the only effect on old-base-relevant content is to redescribe what was already determined, then the move is conservative on the old base.
    \item Otherwise the construction depends essentially on ambient categorical data or on a genuinely new witness object $U(A)$, and it is therefore external to the original base.
\end{enumerate}
In particular, existence by universal property does not by itself furnish new internal operator power on $A$ alone.
\end{proposition}

\begin{proof}
A universal property is never posed with respect to $A$ alone. It quantifies over a wider ambient category, over comparison morphisms, and over the uniqueness or mediating property that singles out the witness object. If all of that structure is already fixed tightly enough that the only resulting statement about the original base is a reformulation of information already determined there, then the move is conservative.

In the substantive case, however, the universal witness is either a genuinely new object or depends on comparison data not fixed by the original base alone. Then the construction is external from the standpoint of that base. Uniqueness up to unique isomorphism does not change the diagnosis. It shows that the witness is canonical relative to the ambient framework, not that it is internally generated by the old base without that framework.
\end{proof}

\begin{example}[Free and adjoint constructions]
The free group on a set, the product of a diagram, or an adjoint image of an object may all be canonical up to the appropriate universal equivalence. But that canonicity is relative to an ambient categorical setting. It does not turn the resulting witness object into an operator internal to the original base object alone.
\end{example}

\begin{theorem}[Flagship case-study closure]\label{thm:case-study-closure}
Each of the five flagship families in this section is governed by the same fixed-base dichotomy. A purported strengthening is either conservative on old-base-relevant content because the load-bearing datum is already internally fixed, or else it depends on external data, witnesses, or ambient extension and is not internal to the original base. In the quotient and completion cases, the move may fail carrier preservation outright by changing the object under discussion.
\end{theorem}

\begin{proof}
Combine Propositions~\ref{prop:case-canonicalization}, \ref{prop:case-completion}, \ref{prop:case-model-theory}, \ref{prop:case-forcing}, and \ref{prop:case-universal}. The point of the section is precisely that the five families are not five unrelated exceptions. They are five recurring sites at which the same internal/external classification reappears.
\end{proof}

\begin{remark}
This section closes the family-by-family part of the argument. The only remaining loophole is interaction: perhaps composition, iteration, diagonalization, or representation change among already classified families produces genuinely new internal operator power. \Cref{sec:no-emergence} addresses exactly that possibility.
\end{remark}

\section{No Emergent Internal Power Under Combination}\label{sec:no-emergence}

After \Cref{thm:case-study-closure}, the remaining loophole is interaction rather than family-by-family classification. One can still ask whether combinations of already classified moves---composition, iteration, diagonal packaging, or representation change---produce a genuinely new internal mechanism. The claim of this section is that they do not. Interaction can increase mathematical reach, but it does not create a third status beyond conservative readback and external dependence.

\begin{definition}[Interaction-generated closure]\label{def:interaction-closure}
Let $\mathcal{F}_0$ be the class of purported carrier-preserving operators on a fixed base structure $\B$ already classified by \Cref{prop:five-disposition,thm:case-study-closure}. The \emph{interaction-generated closure} $\langle \mathcal{F}_0\rangle$ is the smallest class of old-base-relevant procedures containing $\mathcal{F}_0$ and closed under:
\begin{enumerate}[label=(\roman*)]
    \item finite composition;
    \item finite iteration;
    \item admissibility-preserving coding, diagonal packaging, or fixed-point presentation of an already classified procedure; and
    \item admissibility-preserving representation change, intertranslation, or temporary foundation change followed by readback to old-base-relevant content.
\end{enumerate}
\end{definition}

\subsection{Composition and iteration}

\begin{lemma}[Conservative closure under composition]\label{lem:cons-composition}
If $F$ and $G$ each act conservatively on the old-base-relevant content of a fixed base structure $\B$, then $G\circ F$ acts conservatively on that same content.
\end{lemma}

\begin{proof}
An operator that acts conservatively on old-base-relevant content adds no new values for old operations on old tuples and no new old-base consequences. Composing two such operators preserves that status.
\end{proof}

\begin{theorem}[Composition and iteration do not create a third case]\label{thm:composition-iteration}
Let $F_1,\dots,F_n$ be procedures on a fixed base structure $\B$, each already classified as either conservative on old-base-relevant content or external to $\B$. Then the composite
\[
F_n\circ\cdots\circ F_1
\]
has old-base-relevant effect of exactly one of the following two kinds.
\begin{enumerate}[label=(\alph*)]
    \item It is conservative on the old base.
    \item It depends essentially on an external parameter, witness, evaluator, or ambient enlargement and is therefore external to the old base.
\end{enumerate}
In particular, finite iteration of a conservative operator remains conservative, and finite iteration of an externally dependent operator does not become internal merely by repetition.
\end{theorem}

\begin{proof}
Proceed by induction on $n$. The case $n=1$ is the standing classification. If the $(n-1)$-stage composite is conservative and $F_n$ acts conservatively on old-base-relevant content, then the full composite is conservative by \Cref{lem:cons-composition}. If the $(n-1)$-stage composite is already external and its old-base-relevant effect still depends on the load-bearing external datum, then no later stage can make that datum internally fixed; the full composite remains external by \Cref{thm:auxiliary-dichotomy,cor:no-internalization-by-relabeling}.

The only other possibility is that one or more earlier stages use an external detour but the final readback to old-base-relevant content no longer depends on the external witness except as a proving environment. Then the effect of the total composite is conservative rather than internally generative. This is the fixed-base status of many forcing, completion, and monster-model arguments when only an invariant statement about the original base is read back.
\end{proof}

\begin{corollary}[No internal bootstrapping by repetition]\label{cor:no-bootstrapping}
Repeated application of a fixed procedure does not convert external dependence into internal operator power on the original base. Any apparent gain from iteration is either conservative stabilization on old-base-relevant content or further dependence on external data.
\end{corollary}

\begin{proof}
Apply \Cref{thm:composition-iteration} to the iterates $F,F^2,\dots,F^n$.
\end{proof}

\subsection{Diagonalization and fixed-point detours}

Diagonalization and fixed-point arguments are the most plausible remaining loophole because self-reference can look like a source of fresh internal strength. The fixed-base question, however, is the familiar one: which coding, substitution, evaluation, or semantic resources are doing the work?

\begin{proposition}[Diagonalization and fixed-point packaging]\label{prop:diagonalization}
Let $D$ be a purported carrier-preserving operator on a fixed base structure $\B$ whose presentation proceeds by coding, self-application, diagonalization, or fixed-point packaging. Then exactly one of the following holds.
\begin{enumerate}[label=(\alph*)]
    \item The relevant coding, substitution, and readback apparatus are already internally fixed by $\B$. In that case adjoining $D$ amounts only to a conservative definitional expansion or admissible repackaging.
    \item The construction depends essentially on an external satisfaction device, evaluation rule, witness structure, or ambient semantic framework not internally fixed by $\B$. In that case $D$ is external to $\B$.
\end{enumerate}
Therefore diagonalization and fixed-point methods do not create new internal operator power on a fixed base merely by virtue of self-reference.
\end{proposition}

\begin{proof}
If the relevant coding and substitution maps are already fixed by the base, then the purported diagonal operator only repackages material already available in the original regime; in the language of \Cref{sec:exhaustion}, it is conservative definitional expansion or admissible redescription. If stronger evaluative apparatus is required, then the operator depends on auxiliary data not fixed by the old base and is external by \Cref{thm:auxiliary-dichotomy}. Finally, self-reference does not bypass \Cref{thm:no-internal-totalization}: any claim that diagonalization makes an old partial operation newly defined on an old tuple still changes the old graph on old data.
\end{proof}

\subsection{Representation change and foundation swap}

A second common reply is that one can change presentation, recode the subject matter, or move temporarily to a different foundational background and thereby expose a new operator. The paper's answer is again classificatory rather than eliminative: representation change matters in practice, but its fixed-base status is still either conservative or external \citep{ChangKeisler1990,BarrettHalvorson2016,MacLane1998,Simpson2009}.

\begin{proposition}[Representation change and foundation swap]\label{prop:representation-change}
Let $T$ be a procedure that changes notation, coding, signature, presentation, or ambient foundational background and then reads results back to a fixed original base $\B$. Then exactly one of the following holds.
\begin{enumerate}[label=(\alph*)]
    \item The translation preserves and reflects the old-base-relevant content of $\B$. Then the procedure is conservative on that content.
    \item The translation changes the admissibility regime, the permitted witnesses, the carrier under discussion, or the ambient existence assumptions. Then it is an explicit scope change and not an internal strengthening of the original base.
\end{enumerate}
Hence changing representation or foundation does not by itself create new internal operator power on the original base.
\end{proposition}

\begin{proof}
If the translation preserves and reflects old-base-relevant content, then the move is a change of presentation. It may simplify proofs or clarify dependencies, but it does not add new values for old operations on old tuples and does not alter the old-base consequences. Relative to the original base, it is conservative.

If instead the move changes what counts as admissible, what witnesses may be invoked, what objects are available, or what carrier is under discussion, then the procedure is no longer measured against the original regime alone. It is an open change of scope of the kind already isolated in \Cref{prop:five-disposition}. Reading a conclusion back to the original base may still yield a conservative theorem about that base, but the route is external rather than internally generative.
\end{proof}

\begin{theorem}[No-emergent-family theorem]\label{thm:no-emergent-family}
Let $\B$ be a fixed base structure, and let $F\in\langle\mathcal{F}_0\rangle$ be any old-base-relevant procedure obtained from the already classified families by finite composition, iteration, diagonal packaging, or representation/foundation detour with readback. Then the effect of $F$ on the old base is either conservative or external. In particular, interaction among the classified families does not generate a new source of internal operator power on $\B$.
\end{theorem}

\begin{proof}
Every generator in $\mathcal{F}_0$ is already classified by \Cref{prop:five-disposition,thm:case-study-closure}. \Cref{thm:composition-iteration} handles finite composition and iteration, \Cref{prop:diagonalization} handles diagonal and fixed-point packaging, and \Cref{prop:representation-change} handles coding, translation, and foundation change. Hence every element of the closure $\langle\mathcal{F}_0\rangle$ is again either conservative on old-base-relevant content or external to the original base.
\end{proof}

\begin{corollary}[No internal re-entry by detour]\label{cor:no-reentry}
No procedure that leaves a fixed base for external work and later returns with a readback thereby becomes a new internal operator on that base unless the readback is conservative on the old base.
\end{corollary}

\begin{proof}
Immediate from \Cref{thm:no-emergent-family}. The detour may be mathematically essential, but its fixed-base status is either external or conservative.
\end{proof}

\section{Main Theorem and Foundational Consequence}\label{sec:main-theorem}

The preceding sections now provide the full chain of argument. \Cref{thm:no-internal-totalization} blocks direct enlargement of old primitive graphs on old tuples. \Cref{thm:auxiliary-dichotomy} classifies purported operators by the status of their load-bearing auxiliary data. \Cref{thm:exhaustion-theorem} and Proposition~\ref{prop:five-disposition} show that the standard apparent strengthening mechanisms are exhausted and disposed of. \Cref{thm:case-study-closure} shows that the same dichotomy governs the flagship families. \Cref{thm:no-emergent-family} then closes the interaction loophole. What remains is to state the resulting theorem in one place.

\begin{theorem}[No new internal operator power on a fixed base]\label{thm:main-no-new-power}
Let $\B=(X,\Omega,\A)$ be a nondegenerate base structure. Let $F$ be a purported carrier-preserving operator on $\B$ belonging to the interaction-generated closure of the strengthening mechanisms classified in \Cref{sec:exhaustion,sec:no-emergence}. Then exactly one of the following holds.
\begin{enumerate}[label=(\alph*)]
    \item $F$ acts conservatively on the old-base-relevant content of $\B$.
    \item $F$ depends essentially on external parameters, witnesses, evaluators, or ambient enlargement and is therefore not internal to $\B$.
\end{enumerate}
Consequently:
\begin{enumerate}[label=(\roman*)]
    \item no such operator totalizes a genuinely partial old primitive on old tuples;
    \item no composition, iteration, diagonal packaging, or representation change among the classified families yields a new carrier-preserving source of admissible operator power on the original base.
\end{enumerate}
\end{theorem}

\begin{proof}
By \Cref{thm:exhaustion-theorem}, every apparent internal strengthening of the fixed base falls under one of the mechanisms in \Cref{def:five-mechanisms}. Proposition~\ref{prop:five-disposition} shows that each such mechanism is either blocked or else reveals that no genuine internal strengthening has occurred. \Cref{thm:case-study-closure} applies that classification to the flagship mathematical families treated in \Cref{sec:case-studies}. Finally, \Cref{thm:no-emergent-family} shows that the same classification is stable under the natural interaction operations by which those families might be recombined.

Clause (a) collects the conservative cases, and clause (b) collects the externally dependent cases. The first consequence is just the rigidity result \Cref{thm:no-internal-totalization} carried through the later sections. The second consequence is the content of \Cref{thm:no-emergent-family}. No third possibility remains.
\end{proof}

\begin{corollary}[No internal re-entry principle]\label{cor:no-internal-reentry-principle}
External methods may establish theorems, transfer principles, canonical readbacks, and invariance claims about a fixed base. What they do not do is return as new internal operator power on that base unless the return map is conservative on old-base-relevant content.
\end{corollary}

\begin{proof}
Immediate from \Cref{thm:main-no-new-power,cor:no-reentry}.
\end{proof}

\subsection*{Foundational consequence}

The intended consequence is therefore exact rather than rhetorical. The paper does not say that external methods are dispensable, or that their outputs are merely verbal, or that all mathematically important distinctions collapse to definitional equivalence. It says something narrower and, for that reason, more defensible. Once a base structure and its admissibility regime are held fixed, there is no additional internal operator power to be extracted from those methods except by conservative readback.

That conclusion supports a foundational ``no internal re-entry'' principle. Completions, monster models, forcing extensions, universal-property arguments, and other ambient methods remain legitimate and often indispensable. They can reveal structure, establish invariance, prove existence, and reorganize mathematical understanding. But when their conclusions are transported back to the original base, the fixed-base status of that transport is exhausted by the two cases isolated above: either the transport is conservative on the old base, or else it still depends on the ambient extension and is not an internal operator on the original base.

\section{Conclusion}

The paper proves a fixed-base boundary theorem. Relative to a base structure equipped with an admissibility regime, purported new operator power is either conservative on old-base-relevant content or else depends on external witnesses, evaluators, or ambient enlargement. In particular, no carrier-preserving extension totalizes a genuinely partial old primitive on old tuples, and no recombination of the standard detours considered here yields a third internal case.

That conclusion leaves full room for the mathematical strength of completion, forcing, saturation, monster-model arguments, and universal constructions. Their role is not denied but reclassified: they are external methods, and when their results are brought back to the original base the readback is either conservative or still dependent on the external setting. Appendix~\ref{app:extended-cases} records longer reductions intended to make that classification transparent family by family.


\appendix
\section{Extended Family Proofs}\label{app:extended-cases}

The main text isolates, for each flagship family, the load-bearing datum or ambient move and then applies the formal spine developed in \Cref{sec:rigidity,sec:auxiliary-data,sec:exhaustion}. This appendix records longer reductions and mixed-family examples. The pattern is uniform: first determine whether the maneuver changes the carrier under discussion or imports a selector, witness, evaluator, or comparison map; then apply \Cref{thm:no-internal-totalization,thm:auxiliary-dichotomy,prop:five-disposition,thm:composition-iteration}. The appendix introduces no new theorem family; it makes explicit how the main classification works inside recurring subcases.

\subsection{Selection rules and quotients}

\begin{proposition}[Representative rules factor through auxiliary data]\label{prop:app-representatives}
Let $D\subseteq X$ be admissible in a base structure $\B=(X,\Omega,\A)$, and let $\sim$ be an admissibility-invariant equivalence relation on $D$. Suppose a purported carrier-preserving normalization $N:D\to D$ is constant on $\sim$-classes and satisfies $N(x)\sim x$ for all $x\in D$. Then the fixed-base status of $N$ is determined entirely by the status of the representative rule
\[
[x]_{\sim}\longmapsto N(x).
\]
If that rule is internally fixed, $N$ is conservative on old-base-relevant content; if not, $N$ is external. If one passes instead to the quotient $D/{\sim}$, the move is not carrier-preserving.
\end{proposition}

\begin{proof}
Because $N$ is constant on classes, its effect factors through the quotient map $q:D\to D/{\sim}$ together with a choice of section on the relevant quotient classes. The section is the load-bearing datum. When the section is internally fixed by the base---for example by a normalization convention, order, or coding already licensed in the regime---the output of $N$ does not add new old-base content beyond what was already determined by the class of the input. The move is therefore conservative by \Cref{thm:auxiliary-dichotomy}.

When the section is not internally fixed, the normalization depends on an imported order, ranking, code, or tie-breaking policy. In that case the operator depends on an external selector and is not internal to $\B$. Finally, if the intended maneuver is not selection of a representative in $D$ but replacement of each element by its equivalence class, then the effective object is $D/{\sim}$ rather than $D$ itself. That is exactly the carrier-change case isolated in \Cref{prop:five-disposition}.
\end{proof}

\begin{remark}
This decomposition clarifies what is and is not proved in the canonicalization case. The paper does not deny that a representative rule may be mathematically canonical. It says only that, relative to a fixed base, such canonicity must be carried either by an internally fixed normalization policy or by additional structure that is external to the original base.
\end{remark}

\subsection{Completion arguments as enlargement plus readback}

\begin{proposition}[Completion decomposes into enlargement and readback]\label{prop:app-completion-readback}
Let $A$ be an admissible object in a fixed base structure, let $i:A\hookrightarrow \widehat{A}$ be a completion-style embedding into a richer object, and let $R$ be a procedure that reads old-base-relevant consequences about $A$ from work carried out in $\widehat{A}$. Then exactly one of the following occurs.
\begin{enumerate}[label=(\alph*)]
    \item The substantive output concerns new points, ideal elements, or totalized operations in $\widehat{A}$. Then the move is external by carrier change.
    \item The output concerns only the old base $A$, and the readback map $R$ is already internally fixed. Then the move is conservative on old-base-relevant content.
    \item The output concerns only the old base $A$, but the readback depends on an evaluator, boundary convention, embedding policy, or comparison rule not internally fixed by $A$. Then the move is external.
\end{enumerate}
\end{proposition}

\begin{proof}
The completion step itself is a passage to a richer object. That already settles case (a). For cases (b) and (c), the relevant question is no longer whether $\widehat{A}$ exists but how conclusions about $A$ are extracted from reasoning in $\widehat{A}$. That extraction is carried by the readback rule $R$ together with the embedding $i$. If those data are internally fixed relative to the original base, the completion has no old-base effect beyond a conservative reformulation or justification of information already determined there. If not, the completion remains external because the old-base consequence still depends on imported evaluative structure.
\end{proof}

\begin{example}[Cauchy-sequence presentations]
In a metric completion one may represent new points by Cauchy sequences of old points, by Dedekind-style cuts, or by another equivalent completion device. Uniqueness up to isometry controls the richer ambient object. It does not change the fixed-base diagnosis. The passage to the completion is external; only an internally fixed readback to the old metric space is conservative.
\end{example}

\subsection{Compactness and saturation as witness import}

\begin{proposition}[Witness-import schema for model-theoretic arguments]\label{prop:app-model-witness}
Let $M$ be a structure, let $a\in M$ be old parameters, and suppose a purported operator on $M$ is justified by passing to a larger structure $N\succeq M$, by realizing a type over $M$, or by invoking a compactness-based existence theorem \citep{ChangKeisler1990,Hodges1993,Marker2002}. Then the fixed-base effect is classified as follows.
\begin{enumerate}[label=(\alph*)]
    \item If the relevant witness or realization is already definably fixed in $M$, the move is conservative.
    \item If the argument requires a new element of $N$, a witness model, or an existence theorem whose witness is not part of the original base, the move is external.
    \item If the larger structure is used only as a detour to prove an invariant statement about $M$, the readback is conservative even though the proving route is external.
\end{enumerate}
\end{proposition}

\begin{proof}
A saturation or compactness argument typically has two layers: an ambient witness layer and a final statement about the original parameters. If the witness is already present and definably fixed in $M$, then the ambient move adds no new old-base content; it only packages what was already there. If the witness lies outside $M$ or exists only by appeal to a separate existence theorem, then the operator is not internal to $M$ because its effect depends on imported structure.

The remaining case is common in practice: one passes to a saturated or elementary extension only to establish a theorem about the original model. Then the ambient structure is still external as a proving environment, but the effect on old-base-relevant content is conservative. This is exactly the distinction drawn in \Cref{thm:composition-iteration}: an external detour does not become a new internal operator merely because the final sentence is stated in the old language.
\end{proof}

\begin{example}[Realizing a type]
Suppose one moves from $M$ to a saturated extension in order to realize a type over parameters from $M$ and then uses that realization to prove a statement solely about those original parameters. The realized element is not thereby internal to $M$. What is internal, if anything, is only the final old-language consequence that no longer depends on the realizing witness.
\end{example}

\subsection{Forcing as extension plus absoluteness readback}

\begin{proposition}[Forcing decomposes into witness import and readback]\label{prop:app-forcing-decomposition}
Let $M$ be a fixed set-theoretic base, and suppose a purported operator on $M$ is presented via forcing \citep{Jech2003,Kanamori2008}. Then the analysis separates into two stages.
\begin{enumerate}[label=(\alph*)]
    \item The passage from $M$ to a forcing extension $M[G]$ is external, since the generic object $G$ is not internally fixed by $M$.
    \item Any later conclusion about old parameters from $M$ is either a conservative readback, when obtained by absoluteness or invariance, or else still depends on the external extension and remains external.
\end{enumerate}
Consequently forcing never supplies a third fixed-base status beyond external extension and conservative readback.
\end{proposition}

\begin{proof}
The generic filter or generic real is the load-bearing witness of the forcing argument. Because it is not already part of the original base, the move to $M[G]$ is external. That settles the generative status of the forcing extension itself.

What can still vary is the status of the conclusion one transports back to $M$. If the transported statement concerns only old parameters and is secured by absoluteness, preservation, or an invariance argument, then the effect on old-base-relevant content is conservative: the forcing extension served as an external proving environment. If, however, the purported output still refers to the generic object, to the richer extension, or to data determined only there, then the dependence remains external. The same diagnosis applies to Boolean-valued or algebraic presentations of forcing: these may redescribe the route, but they do not internalize the generic witness.
\end{proof}

\begin{remark}
This is the clearest example of the paper's ``no internal re-entry'' theme. One may leave the base, reason in a richer universe, and return with a theorem about the base. What does not happen is that the generic object thereby becomes an internally generated operator on the original universe.
\end{remark}

\subsection{Universal properties and ambient comparison data}

\begin{proposition}[Universal witnesses depend on ambient comparison data]\label{prop:app-universal-comparison}
Let $A$ be an object presented by a fixed base structure, and suppose an associated construction is specified by a universal property in an ambient category $\mathcal{C}$ \citep{MacLane1998}. Then the fixed-base status of the construction is carried by the ambient comparison data: the class of comparison maps, the universal condition, and the witness object satisfying it. If those data are used only to justify a conservative reformulation of old-base content, the effect is conservative; otherwise the move is external.
\end{proposition}

\begin{proof}
A universal property is not posed over $A$ alone. It quantifies over a wider ambient setting and singles out a witness object by its behavior relative to all comparison maps of the specified sort. That witness is therefore not generated by the original base object in isolation. If the only old-base consequence is an invariant reformulation already determined there, the effect is conservative. In the substantive case, however, the universal witness or the comparison condition does real mathematical work and so functions as external ambient data.
\end{proof}

\begin{example}[Free objects]
The free group on a set, a product of a diagram, or a left-adjoint image may be canonical up to unique isomorphism, but that canonicity is relative to the chosen ambient category and class of comparison morphisms. It does not convert the witness object into an operator internal to the original object alone.
\end{example}

\subsection{Mixed-family worked examples}

\begin{proposition}[Mixed detours inherit the same classification]\label{prop:app-mixed}
Let $F$ be an old-base-relevant procedure on a fixed base obtained by finitely many stages, where each stage is one of the family moves analyzed in \Cref{sec:case-studies}. Then the total effect of $F$ on the original base is again either conservative or external.
\end{proposition}

\begin{proof}
This is an instance of \Cref{thm:composition-iteration}. What matters for appendix purposes is how the reduction looks in concrete mixed cases.
\end{proof}

\begin{example}[Completion followed by canonicalization]
One first passes from an old object to a completion, then chooses a ``canonical'' representative or normal form there, and finally reads a statement back to the old base. The completion stage is external. If the representative rule used in the completed object is not internally fixed even relative to that ambient setting, a second external dependence has been added. If the final readback to the old base forgets those witnesses and states only an invariant fact about the original object, the net old-base effect is conservative; otherwise the entire procedure remains external.
\end{example}

\begin{example}[Forcing followed by absoluteness]
One passes from $M$ to $M[G]$, proves a statement about old parameters, and then invokes absoluteness to transfer the statement back to $M$. The first stage imports the generic witness and is external. The final transfer is conservative on old-base-relevant content precisely because the transported statement no longer depends on $G$ as part of its content. This is the template behind many fixed-base readback arguments.
\end{example}

\begin{example}[Monster-model detour followed by quotient or canonical parameter]
One works in a monster model to compare types or to define a quotient-like invariant and then states a theorem about the original small model. The monster model is an ambient enlargement, and any quotient passage or representative choice must still be classified by the corresponding family analysis. No matter how the stages are combined, the old-base effect is either a conservative theorem about the small model or else an external dependence on the larger ambient setting.
\end{example}

\section{Glossary of Terminology}\label{app:glossary}

\begin{center}
\begin{tabular}{ll}
\toprule
Current Term & Submission Term \\
\midrule
standing-bearing carrier & fixed base structure / fixed carrier \\
same-scope operator & carrier-preserving / internal operator \\
bookkeeping-equivalent & conservative redescription / conservative extension \\
standing-fixed datum & internally fixed auxiliary datum \\
imported datum & external parameter / external witness \\
AMetric boundary & external completion boundary / evaluator \\
\bottomrule
\end{tabular}
\end{center}

\section{Dependency Map and Proof Chart}\label{app:dependency-map}

This appendix isolates the formal spine of the manuscript. Appendix~\ref{app:extended-cases} supplies expanded applications of the spine but introduces no additional premises.

The main dependency chain is:
\begin{enumerate}[label=D\arabic*., leftmargin=3em]
    \item \Cref{lem:old-graph-rigidity,lem:domain-fixity,thm:no-internal-totalization}: rigidity of carrier-preserving extension;
    \item \Cref{thm:auxiliary-dichotomy}: internal/external auxiliary-data classification;
    \item \Cref{thm:exhaustion-theorem,prop:five-disposition,cor:closure-totalization}: exhaustion of apparent strengthening mechanisms;
    \item the case-study propositions in \Cref{sec:case-studies} and the expanded instances in Appendix~\ref{app:extended-cases};
    \item \Cref{thm:composition-iteration,prop:diagonalization,prop:representation-change,thm:no-emergent-family}: closure under recombination;
    \item \Cref{thm:main-no-new-power}: final synthesis.
\end{enumerate}


{\small
\begin{thebibliography}{99}

\bibitem[Barrett and Halvorson(2016)]{BarrettHalvorson2016}
Thomas William Barrett and Hans Halvorson.
\newblock Morita equivalence.
\newblock \emph{The Review of Symbolic Logic}, 9(3):556--582, 2016.

\bibitem[Barwise(1975)]{Barwise1975}
Jon Barwise.
\newblock \emph{Admissible Sets and Structures}.
\newblock Perspectives in Mathematical Logic. Springer-Verlag, Berlin, 1975.

\bibitem[Burmeister(1986)]{Burmeister1986}
Peter Burmeister.
\newblock \emph{A Model Theoretic Oriented Approach to Partial Algebras: Introduction to Theory and Application of Partial Algebras -- Part I}.
\newblock Mathematical Research 32. Akademie-Verlag, Berlin, 1986.

\bibitem[Chang and Keisler(1990)]{ChangKeisler1990}
C.~C. Chang and H.~Jerome Keisler.
\newblock \emph{Model Theory}.
\newblock North-Holland, Amsterdam, third edition, 1990.

\bibitem[Hodges(1993)]{Hodges1993}
Wilfrid Hodges.
\newblock \emph{Model Theory}.
\newblock Encyclopedia of Mathematics and its Applications 42. Cambridge University Press, Cambridge, 1993.

\bibitem[Jech(2003)]{Jech2003}
Thomas Jech.
\newblock \emph{Set Theory: The Third Millennium Edition, Revised and Expanded}.
\newblock Springer Monographs in Mathematics. Springer, Berlin, 2003.

\bibitem[Kanamori(2008)]{Kanamori2008}
Akihiro Kanamori.
\newblock The mathematical development of set theory from {Cantor} to {Cohen}.
\newblock \emph{The Bulletin of Symbolic Logic}, 14(1):1--71, 2008.

\bibitem[Mac Lane(1998)]{MacLane1998}
Saunders Mac Lane.
\newblock \emph{Categories for the Working Mathematician}.
\newblock Graduate Texts in Mathematics 5. Springer, New York, second edition, 1998.

\bibitem[Marker(2002)]{Marker2002}
David Marker.
\newblock \emph{Model Theory: An Introduction}.
\newblock Graduate Texts in Mathematics 217. Springer, New York, 2002.

\bibitem[Simpson(2009)]{Simpson2009}
Stephen~G. Simpson.
\newblock \emph{Subsystems of Second Order Arithmetic}.
\newblock Perspectives in Logic. Cambridge University Press, Cambridge, second edition, 2009.

\end{thebibliography}

}\normalsize


\end{document}
